\section{Step 3: verify the classical benchmark (current)}
\label{sec:stage3}

\subsection*{Main equation and assumptions}

This stage specializes the exact-liquidation equation from Step 2:
\[
 \boxed{\ddot{\mathbf x}=C\mathbf x,\qquad
 \mathbf x(0)=(3,1)^\top,\quad\mathbf x(1)=0,\qquad
 \mathbf v=-\dot{\mathbf x}.}
\]
It computes the trajectory, checks whether its signed optimum is
sales-only, and derives quantities for numerical comparison.

\begin{enumerate}
\item Use the constant quadratic model of Step 2 with $H=I_2$,
$\lambda=T=1$. All quantities in this synthetic instance are
dimensionless. There is no calibration to observed market data.
\item Impose the hard endpoint $\mathbf x(1)=0$, rather than a finite
terminal penalty. The matrix $P(t)$ therefore has a singular terminal
limit; no finite value of $P(1)$ is substituted below.
\item First optimize over signed rates with no inventory boundaries or
rate caps. Sales-only feasibility is a conclusion to prove for
$\mathbf X=(3,1)^\top$, not an assumption used to solve the ODE.
\item The general ratio test below requires equal impacts, equal price
variances, positive correlation $0<\rho<1$, and positive initial holdings
$X_1=\theta X_2$ with $\theta>1$. These restrictions make the
symmetric/antisymmetric modes and the terminal-ratio proof valid.
\end{enumerate}

\subsection*{Diagonalize the trajectory equation}

Use the synthetic, dimensionless instance
\[
 H=I_2,\qquad C=\begin{pmatrix}1&0.4\\0.4&1\end{pmatrix},\qquad
 \lambda=T=1,\qquad\mathbf X=(3,1)^\top,\qquad\mathbf x(1)=0.
\]
The orthonormal covariance eigenvectors are
$u_+=(1,1)^\top/\sqrt2$ and $u_-=(1,-1)^\top/\sqrt2$ with
\[
 \mu_+=1.4,\quad\mu_-=0.6,\qquad
 \kappa_+=\sqrt{1.4}=1.1832159566,\quad
 \kappa_-=\sqrt{0.6}=0.7745966692.
\]
Thus $C\succ0$, $\eta_j>0$, and both rates are real. Set $s=1-t$ and
use $U=(u_+,u_-)$, so $U^\top C U=\operatorname{diag}(1.4,0.6)$.
Writing $\mathbf y=U^\top\mathbf x$ gives
\[
 \ddot y_\pm=\kappa_\pm^2y_\pm,\qquad
 y_+(0)=\frac{3+1}{\sqrt2}=2\sqrt2,\qquad
 y_-(0)=\frac{3-1}{\sqrt2}=\sqrt2,\qquad y_\pm(1)=0.
\]
The endpoint solution from Step 2 is
$y_\pm(t)=y_\pm(0)\sinh[\kappa_\pm(1-t)]/\sinh\kappa_\pm$.
Define
\[
 f_\pm(s)=\frac{\sinh(\kappa_\pm s)}{\sinh\kappa_\pm},\qquad
 a_\pm(s)=\frac{\kappa_\pm\cosh(\kappa_\pm s)}{\sinh\kappa_\pm}.
\]
Since $\mathbf x=U\mathbf y$, one has
$x_1=(y_++y_-)/\sqrt2$ and $x_2=(y_+-y_-)/\sqrt2$.
Also $ds/dt=-1$, so $df_\pm(1-t)/dt=-a_\pm(1-t)$.
These two substitutions give the trajectory and selling rates
\begin{equation}
 x_1=2f_++f_-,\quad x_2=2f_+-f_-,\qquad
 v_1=2a_++a_-,\quad v_2=2a_+-a_-.
 \label{eq:stage3-rates}
\end{equation}
They give $\mathbf x(0)=(3,1)^\top$ and $\mathbf x(1)=0$.

\subsection*{Prove the rates are sales-only for all times}

For the equal-impact, equal-volatility family with $0<\rho<1$ and
$X_1=\theta X_2$, $\theta>1$, $X_2>0$, the potentially negative rate is
\[
 v_2=\frac{X_2}{2}\big((\theta+1)a_+-(\theta-1)a_-\big).
\]
The relative rate $\mathcal R(s)=a_-(s)/a_+(s)$ satisfies
\[
 \frac{d}{ds}\log \mathcal R(s)
 =\kappa_-\tanh(\kappa_-s)-\kappa_+\tanh(\kappa_+s)\le0,
\]
because $k\mapsto k\tanh(ks)$ increases for $k>0$, $s\ge0$.
Indeed its derivative is
$\tanh(ks)+ks\operatorname{sech}^2(ks)\ge0$.
The rate $v_1$ is positive because both modal coefficients are positive.
For $v_2$, divide by the positive factor
$X_2(\theta-1)a_+/2$ to obtain
\[
 v_2\ge0\quad\Longleftrightarrow\quad
 \frac{\theta+1}{\theta-1}\ge\frac{a_-}{a_+}=\mathcal R(s).
\]
Hence $\mathcal R(s)\le \mathcal R(0)$: checking the terminal-time ratio is sufficient and
necessary for this particular family, rather than an assumption about all
multi-asset problems. Specifically,
\begin{equation}
 \frac{\theta+1}{\theta-1}\ge
 \mathcal R(0)=\frac{\kappa_-\sinh\kappa_+}{\kappa_+\sinh\kappa_-}
 =1.133429601920,
 \qquad
 \theta\le\frac{\mathcal R(0)+1}{\mathcal R(0)-1}=15.989177597983.
 \label{eq:stage3-bound}
\end{equation}
The choice $\theta=3$ satisfies this bound. More explicitly, since
$a_+(s)\ge a_+(0)$ and $2-\mathcal R(s)\ge2-\mathcal R(0)>0$,
\[
 v_2(t)\ge a_+(0)(2-\mathcal R(0))=0.6931333103,
 \qquad v_1(t)\ge2a_+(0)+a_-(0)=2.5062988967.
\]
Both minima occur at $t=1$; these are one-sided terminal limits of the rates.
Thus both inventories stay nonnegative and decrease to zero for the entire
interval. A time-grid check alone would not establish this all-time claim.

\subsection*{Derive the optimal cost and initial rate}

The initial modal inventories are $y_+(0)=2\sqrt2$ and
$y_-(0)=\sqrt2$. Orthogonality of $U$ gives
\[
 J^*=\sum_{i\in\{+,-\}}\int_0^1
 (\dot y_i^2+\kappa_i^2y_i^2)\,dt.
\]
Use $\ddot y_i=\kappa_i^2y_i$ to recognize a boundary derivative:
\[
 \frac{d}{dt}(y_i\dot y_i)=\dot y_i^2+y_i\ddot y_i
 =\dot y_i^2+\kappa_i^2y_i^2.
\]
The terminal term vanishes because $y_i(1)=0$, and
$\dot y_i(0)=-\kappa_i\coth\kappa_i\,y_i(0)$. Consequently
\begin{equation}
 \boxed{J^*=8\kappa_+\coth\kappa_+
              +2\kappa_-\coth\kappa_-
 =\tfrac12\mathbf X^\top P(0)\mathbf X.}
 \label{eq:stage3-cost}
\end{equation}
The same mode derivatives give
\[
 \mathbf v(0)=
 \begin{pmatrix}
 2\kappa_+\coth\kappa_++\kappa_-\coth\kappa_-\\
 2\kappa_+\coth\kappa_+-\kappa_-\coth\kappa_-
 \end{pmatrix}.
\]
These expressions explain the numerical targets below without a
numerical integration of the trajectory.

\subsection*{Checks to reproduce}

Targets retained from the existing notes are
\[
 \mathbf v(0)=(4.0488436487,1.6639810858)^\top,\qquad
 J^*=\tfrac12\mathbf X^\top P(0)\mathbf X=13.810512031917.
\]
Recompute the cost by direct quadrature and compare it with the Riccati value.
The full numerical check script is absent; the table below remains a set
of targets to reproduce, not completed verification.

For an independent check, minimize
$\sum_{k=1}^N\|\mathbf x_{k-1}-\mathbf x_k\|^2/\tau
+\tau\sum_{k=1}^N\mathbf x_k^\top C\mathbf x_k$ by solving its block
tridiagonal stationarity system. Its Hessian is positive definite, so this
is the unique global discrete minimizer.
Here $\tau=1/N$, $\mathbf x_0=\mathbf X$ and $\mathbf x_N=0$.
Differentiating the discrete objective with respect to an interior
$\mathbf x_k$, $1\le k<N$, gives
\[
 \frac{2}{\tau}(2\mathbf x_k-\mathbf x_{k-1}-\mathbf x_{k+1})
 +2\tau C\mathbf x_k=0,
 \qquad
 (2I+\tau^2 C)\mathbf x_k-\mathbf x_{k-1}-\mathbf x_{k+1}=0.
\]
Thus the diagonal blocks are $2I+\tau^2 C$, the neighboring blocks
are $-I$, and the first right-hand-side block is $\mathbf X$.
For a nonzero variation with zero endpoints, the quadratic form of the
objective's Hessian is
\[
 \frac{2}{\tau}\sum_{k=1}^N
 \|\delta\mathbf x_k-\delta\mathbf x_{k-1}\|^2
 +2\tau\sum_{k=1}^N\delta\mathbf x_k^\top C\delta\mathbf x_k>0.
\]
This proves the discrete uniqueness assertion and specifies the
linear system needed to reproduce the table.
\begin{center}
\begin{tabular}{@{}rrr@{}}
\toprule
$N$ & Maximum inventory error vs. continuum & Discrete optimal cost \\
\midrule
16 & $6.945\times10^{-5}$ & 13.432845715492 \\
64 & $4.361\times10^{-6}$ & 13.714251771181 \\
256 & $2.726\times10^{-7}$ & 13.786331703645 \\
1024 & $1.704\times10^{-8}$ & 13.804459745777 \\
\bottomrule
\end{tabular}
\end{center}
The trajectory error decreases quadratically; the cost has first-order
right-endpoint quadrature error. These are different error measures.
The permanent-impact/spread offsets are not part of these reported costs.

\subsection*{Source, scope and exit criterion}

The trajectory comes from Step 2's specialization of
Bergault--Drissi--Gu\'eant \cite{bergault2022} and its classical
mode solution, also described in Almgren--Chriss
\cite{almgrenchriss2000}, Appendix~A.3. The parameter choice and sales-only
inequality are derived here. The all-time proof applies to the stated
hard-liquidation instance. A finite penalty needs its own feasibility check,
and $H=I$ does not test a noncommuting implementation.

\textbf{Pending.} Review the Eq.~16 mapping, Riccati equation, terminal data
and sales-only proof. Reproduce the cost and convergence table, and compare
the block-exponential finite-penalty formula of Step 2 with direct Riccati integration for a
noncommuting penalty/covariance pair. Record the checks, then begin Step 4.

